\begin{pastexam}[label={px:ntu-109-comm-dsa-24}]{臺大電信所 109 學年度（B）第 24 題}
（單選；每題 2 分，答錯不倒扣）

In this exam, the height of a singleton tree is defined as 0.

An AVL tree of 100 nodes has height at least

(A)~5 \qquad (B)~6 \qquad (C)~7 \qquad (D)~8 \qquad (E)~9
\end{pastexam}

\begin{pastsoln}
答案是 (B)。

本卷規定只有一個節點的樹高度是 $0$，也就是高度以\textbf{邊數}計；資結書以節點數計（單一
節點高度 $1$），換算時減 $1$。

\textbf{下界。} 任何高度 $h$（邊數）的二元樹最多 $h + 1$ 層，第 $i$ 層（根在第 $0$ 層）最多
$2^i$ 個節點，總共最多 $2^{h+1} - 1$ 個。 要放下 $100$ 個節點需要 $2^{h+1} - 1 \ge 100$：
$h = 5$ 時只有 $63$，不夠；$h = 6$ 時 $127 \ge 100$。 所以 $h \ge 6$。

\textbf{做得到。} $100$ 個節點的完全二元樹（前 $6$ 層 $63$ 個填滿，第 $7$ 層由左往右放 $37$
個）高度是 $6$；它每個節點的左右子樹高度最多差 $1$，是合法的 AVL 樹。 所以 AVL 樹的高度最小
就是 $6$。

順便看另一個方向：高度 $h$ 的 AVL 樹最少有 $N(h)$ 個節點，$N(0) = 1$、$N(1) = 2$、$N(h) =
N(h - 1) + N(h - 2) + 1$，依序是 $1, 2, 4, 7, 12, 20, 33, 54, 88, 143$（資結書定理~10.1 的
$N$ 以節點數計高度，數列相同、編號差 $1$）。 $N(8) = 88 \le 100 < N(9) = 143$，所以 $100$ 個
節點的 AVL 樹高度\textbf{最多} $8$，那是 (D) 的來歷，但題目問的是 at least。
\end{pastsoln}
